\documentclass[10pt,a4paper]{amsart}
\usepackage[margin=19mm]{geometry}
\usepackage{amsmath,amssymb,mathtools}
\usepackage[scr=rsfs,cal=euler]{mathalfa}
\usepackage{xfrac}
\usepackage[unicode,hidelinks,bookmarks=false]{hyperref}
\newcommand{\ie}{i.e.\ }
\newcommand{\cf}{cf.\ }
\newcommand{\resp}{respectively\ }
\newcommand{\Subsection}{\subsection}
\providecommand{\nobreakdash}{\nobreak\hbox{-}}
\setlength{\emergencystretch}{2em}
\multlinegap0pt
\usepackage{xcolor}
\usepackage{amssymb}
\usepackage[shortlabels]{enumitem}
\usepackage{tikz}
\usepackage{algorithm}\usepackage{algpseudocode}
\newtheorem{theorem}{Theorem}
\usepackage{cleveref}
\crefname{theorem}{Theorem}{Theorems}
\newcommand{\Deg}{\mathrm{d}}
\newcommand{\Maas}{\mathcal{M}}
\newcommand{\Prob}{\mathscr{P}}
\newcommand{\Vol}{\mathrm{Vol}}
\title{Resistance Distance and \\ Linearized Optimal Transport on Graphs}
\author{Sawyer Jack Robertson\,, Zhengchao Wan\,, and Alexander Cloninger${}^\ast{}$}
\date{Source: arXiv:2404.15261v4 (CC BY 4.0)}
\begin{document}
\maketitle

\noindent\emph{Source and licence.} Resistance Distance and \\ Linearized Optimal Transport on Graphs. Version arXiv:2404.15261v4 (CC BY 4.0), distributed by arXiv under the Creative Commons Attribution 4.0 International licence, http://creativecommons.org/licenses/by/4.0/, as recorded in the arXiv licence metadata for this exact version and shown on the abstract page https://arxiv.org/abs/2404.15261v4. Full licence notice: \texttt{licenses/arxiv-2404.15261-CC-BY-4.0.txt}.\par\smallskip

\noindent\textit{Editorial scope.} Counted unit: the article's theorem thm:linearization-mu (source0.tex lines 331--347). The article's complete original proof of it is delivered with it (source0.tex lines 349--536, one \texttt{proof} environment of 188 lines). No mathematics is written, repaired or re-derived: the statement and the proof are the article's own lines, in the article's own notation, and no retained line is substituted or re-flowed.  Retained uncounted as the source-local context the proof needs: lines 145--149.  Nothing was omitted from any retained span, so the package carries no omission record.  No retained line contains a citation, so the wrapper carries no bibliography and no bibliography entry.  No retained line contains a figure environment, an image-inclusion command or drawing code, so no image is delivered and the package declares no asset dependency.  Editorial formatting only, in addition: an AMS \texttt{amsart} preamble that restates the article's own declarations and macros used by the retained lines, namely the environments \texttt{theorem} on separate counters, together with the standard packages those lines need.  The retained environments print in this document as Theorem 1 in order of appearance, whereas the article prints the same items with the numbering of its own full document.  The complete native source of the article as distributed in the arXiv e-print for this exact version is bundled unchanged as \texttt{sources/158059/source0.tex}.
        \begin{align}\label{eq:defn-maas-distance}
            \Maas (\mu,\nu)^2 &= \inf\Bigl\{ \mathcal{A}((\psi_t), (\rho_t)) \;:\; (\rho_t)_{t\in[0,1]}\in \mathcal{C}^{1}([0, 1];\,\Prob_\star(V)),\; (\psi_t)_{t\in[0,1]}\in \mathcal{B}([0, 1];\,\mathbb{R}^{V}),\notag\\
            &\qquad\qquad \frac{d}{dt}\rho_t(x) + \sum_{y\sim x} w_{xy}\, (\psi_t(y) - \psi_t(x))\, \theta\left(\frac{\rho_t(x)}{\pi(x)}, \frac{\rho_t(y)}{\pi(y)}\right) = 0,\notag\\
            &\qquad\qquad \rho_0 = \mu,\ \rho_1 = \nu \Bigr\}.
        \end{align}

\begin{theorem}\label{thm:linearization-mu}
    Let $G=(V,E, w)$ be a connected graph and let $\pi\in\Prob_\star (V)$ denote the stationary distribution of the simple random walk on $V$, and let $\mu\in\Prob_\star(V)$ be fixed. Define edge weights $w^{(\mu)}=(w^{(\mu)}_{xy})_{\{x,y\}\in E}$ by
        \begin{align*}
            w^{(\mu)}_{xy} := w_{xy}\,\theta\left(\frac{\mu(x)}{\pi(x)},\,\frac{\mu(y)}{\pi(y)}\right),\quad \{x,y\}\in E,
        \end{align*}
    and let $L_{w^{(\mu)}}$ be the corresponding Laplacian. Let $h\in\mathbb{R}^{V}_0$ be fixed with $\|h\|_2\le 1$. Let $\lambda_2(L_{w^{(\mu)}})$ denote the smallest nonzero eigenvalue of $L_{w^{(\mu)}}$, and define constants $0 < C \le C'$ via
        \begin{align}\label{eq:constants-linearization}
            C &:= \left(\min_{x\in V}\mu(x)\right)^{-1},\quad C' := C\,\sqrt{\frac{2\,\Vol_w(G)}{\lambda_2(L_{w^{(\mu)}})}}.
        \end{align}
    Then for any $\varepsilon\in\mathbb{R}$ with $|\varepsilon|$ sufficiently small (depending only on $G$ and $\mu$), it holds
        \begin{align}\label{eq:main-linearization}
            \frac{\varepsilon^2}{1 + C'|\varepsilon|}\, h^\top L_{w^{(\mu)}}^\dagger h
            \le \Maas(\mu,\mu+\varepsilon h)^2
            \le \frac{\varepsilon^2}{1 - C|\varepsilon|}\, h^\top L_{w^{(\mu)}}^\dagger h.
        \end{align}
    In particular, $\Maas(\mu,\mu+\varepsilon h)^2 = |\varepsilon|^2\, h^\top L_{w^{(\mu)}}^\dagger h  + o(|\varepsilon|^2)$.
\end{theorem}

\begin{proof}
    First we will exhibit an admissible pair $(\rho_t), (\psi_t)$ as in~\cref{eq:defn-maas-distance} to obtain the upper bound in~\cref{eq:main-linearization}. To this end let
        \begin{align*}
            \rho_t &= t\, (\mu + \varepsilon h) + (1-t) \, \mu,\quad 0\le t\le 1.
        \end{align*}
    Then if $|\varepsilon|$ is chosen to be small enough, $(\rho_t)\in \mathcal{C}^{1}([0, 1];\,\Prob_\star(V))$. Clearly $\rho_0 = \mu,\, \rho_1=\mu + \varepsilon h$. The flux constraint in~\cref{eq:defn-maas-distance} requires of $\psi_t$ that
        \begin{align}\label{eq:flux-constraint-linear}
            \frac{d}{dt} \rho_t(x) &= \varepsilon h(x) = -\sum_{y\sim x} w_{xy}\, (\psi_t(y) - \psi_t(x))\, \theta\left(\frac{\rho_t(x)}{\pi(x)}, \frac{\rho_t(y)}{\pi(y)}\right),\quad x\in V,\, 0\le t\le 1.
        \end{align} 
    Write
        \begin{align*}
            \widetilde{w}^{(t)}_{xy} &= w_{xy}\,\theta\left(\frac{\rho_t(x)}{\pi(x)}, \frac{\rho_t(y)}{\pi(y)}\right),\quad \{x, y\}\in E,
        \end{align*}
    so that~\cref{eq:flux-constraint-linear} becomes
        \begin{align}\label{eq:flux-constraint-linear-2}
            \varepsilon \, h &= L_{\widetilde{w}^{(t)}} \psi_t,\quad 0\le t\le 1.
        \end{align}
    Note that since $\rho_t\in\Prob_\star (V)$ holds for all $t$ and since $\theta(a, b) > 0$ whenever $a, b > 0$, the weights $\widetilde{w}^{(t)}_{xy}$ remain strictly positive for all $0\le t\le 1$. In particular, the operator $L_{\widetilde{w}^{(t)}}$ remains the Laplacian of a connected weighted graph and is therefore positive definite on the subspace $\mathbb{R}^{V}_0$ which includes $h$. We may therefore pick, for each $0\le t\le 1$, a unique representative $\psi_{t}^\ast\in\mathbb{R}^{V}_0$ satisfying~\cref{eq:flux-constraint-linear-2} and given by
        \begin{align*}
            \psi_t^\ast &= \varepsilon \, L_{\widetilde{w}^{(t)}}^{\dagger} h.
        \end{align*}
    Then, using $L_{\widetilde{w}^{(t)}}^{\dagger}L_{\widetilde{w}^{(t)}} L_{\widetilde{w}^{(t)}}^{\dagger} = L_{\widetilde{w}^{(t)}}^{\dagger}$, the action functional evaluates to
        \begin{align}\label{eq:action-upper-bound-1}
            \mathcal{A}((\psi_t^\ast), (\rho_t)) &= \int_0^{1} (\psi_t^\ast)^\top L_{\widetilde{w}^{(t)}} \psi_t^\ast \, dt = \varepsilon^2 \int_0^{1} h^\top L_{\widetilde{w}^{(t)}}^{\dagger} h \, dt\notag\\
            &= \varepsilon^2 \, h^\top \left( \int_0^{1} L_{\widetilde{w}^{(t)}}^{\dagger} \, dt \right) h.
        \end{align}
    To bound~\cref{eq:action-upper-bound-1} from above we first observe that, for each $x\in V$ and $0\le t\le 1$, one has
        \begin{align*}
            \frac{\rho_t(x)}{\pi(x)} &= \frac{\mu(x) + t \varepsilon h(x)}{\pi(x)} = \frac{\mu(x)}{\pi(x)} + t \varepsilon \frac{h(x)}{\pi(x)}.
        \end{align*}
    Using $|h(x)|\le \|h\|_2\le 1$, we have
        \begin{align*}
            \left|t \varepsilon \frac{h(x)}{\pi(x)}\,\frac{\pi(x)}{\mu(x)}\right| &\leq |\varepsilon| C,
        \end{align*}
    where $C$ is taken as in~\cref{eq:constants-linearization}. Thus it follows that
        \begin{align}\label{eq:pi-mu-bound}
            \left(1-|\varepsilon|C\right)\frac{\mu(x)}{\pi(x)} \le \frac{\rho_t(x)}{\pi(x)} \le \left(1+|\varepsilon|C\right)\frac{\mu(x)}{\pi(x)},\quad x\in V,\,0\le t\le 1.
        \end{align}
    Since $\theta$ is increasing in each argument and homogeneous of degree one,~\cref{eq:pi-mu-bound} yields
        \begin{align*}
            \left(1-|\varepsilon|C\right)\theta\left(\frac{\mu(x)}{\pi(x)},\frac{\mu(y)}{\pi(y)}\right)
            \le \theta\left(\frac{\rho_t(x)}{\pi(x)},\frac{\rho_t(y)}{\pi(y)}\right)
            \le \left(1+|\varepsilon|C\right)\theta\left(\frac{\mu(x)}{\pi(x)},\frac{\mu(y)}{\pi(y)}\right),
        \end{align*}
    and therefore, for each $\{x, y\}\in E$ and $0\le t\le 1$, it holds
        \begin{align*}
            (1 - |\varepsilon| C) \, w^{(\mu)}_{xy} \leq \widetilde{w}^{(t)}_{xy} \leq (1 + |\varepsilon| C) \, w^{(\mu)}_{xy}.
        \end{align*}
    Now let $f\in\mathbb{R}^{V}$ be arbitrary. It follows that
        \begin{align*}
            (1 - |\varepsilon| C) \, f^\top L_{w^{(\mu)}} f &\leq f^\top L_{\widetilde{w}^{(t)}} f \leq (1 + |\varepsilon| C) \, f^\top L_{w^{(\mu)}} f,
        \end{align*}
    and in particular this bound holds for $f\in\mathbb{R}^{V}_0$. Since $L_{w^{(\mu)}}$ and $L_{\widetilde{w}^{(t)}}$ are positive definite on $\mathbb{R}^{V}_0$, the inverse inequality holds for their Moore--Penrose inverses:
        \begin{align}\label{eq:pseudoinverse-bound}
            \frac{1}{1 + |\varepsilon| C} \, f^\top L_{w^{(\mu)}}^{\dagger} f &\leq f^\top L_{\widetilde{w}^{(t)}}^{\dagger} f \leq \frac{1}{1 - |\varepsilon| C} \, f^\top L_{w^{(\mu)}}^{\dagger} f,\quad f\in\mathbb{R}^{V}_0.
        \end{align}
    Since~\cref{eq:pseudoinverse-bound} holds for all $0\le t\le 1$, we may integrate all sides to obtain
        \begin{align*}
            \frac{1}{1 + |\varepsilon| C} \, f^\top L_{w^{(\mu)}}^{\dagger} f &\leq f^\top \left( \int_0^{1} L_{\widetilde{w}^{(t)}}^{\dagger} \, dt \right) f \leq \frac{1}{1 - |\varepsilon| C} \, f^\top L_{w^{(\mu)}}^{\dagger} f,
        \end{align*}
    thus, with $f=h$ and via~\cref{eq:action-upper-bound-1}, one has
        \begin{align}\label{eq:pseudoinverse-bound-2}
            \mathcal{A}((\psi_t^\ast), (\rho_t)) &= \varepsilon^2 \, h^\top \left( \int_0^{1} L_{\widetilde{w}^{(t)}}^{\dagger} \, dt \right) h \leq \frac{\varepsilon^2}{1-|\varepsilon| C} h^\top L_{w^{(\mu)}}^{\dagger} h.
        \end{align}
    
    Next we prove the lower bound. By~\cref{eq:pseudoinverse-bound-2} and taking $|\varepsilon|C\le \frac12$, we have
        \begin{align*}
            \Maas(\mu, \mu+\varepsilon \, h)^2 &\leq 2\varepsilon^2 h^\top L_{w^{(\mu)}}^{\dagger} h 
        \end{align*}
    Since $\|h\|_2\le 1$ and $h\in\mathbb{R}^{V}_0$, we have $h^\top L_{w^{(\mu)}}^{\dagger} h \le \lambda_2(L_{w^{(\mu)}})^{-1}\|h\|_2^2 \le \lambda_2(L_{w^{(\mu)}})^{-1}$, so
        \begin{align*}
            \Maas(\mu, \mu+\varepsilon \, h)^2 &\leq \frac{2}{\lambda_2(L_{w^{(\mu)}})}\varepsilon^2.
        \end{align*}
    Now let $\eta > 0$ be fixed, and let $(\rho_t), (\psi_t)$ be any admissible pair of curves satisfying
        \begin{align*}
            \mathcal{A}((\psi_t), (\rho_t)) &\leq \Maas(\mu, \mu+\varepsilon h)^2 + \eta \leq \frac{2}{\lambda_2(L_{w^{(\mu)}})}\varepsilon^2 + \eta.
        \end{align*}
    We emphasize that $(\rho_t), (\psi_t)$ depend on the choice of $\eta$, but we suppress this in the notation. Then for any $x\in V$ and $0\le t\le 1$, since $(\rho_t)$ is piecewise $\mathcal{C}^1$, one has that
        \begin{align*}
            \rho_t(x) - \rho_0(x) &= \int_0^t \frac{d}{ds} \rho_s(x) \, ds = \int_0^t L_{\widehat{w}^{(s)}} \psi_s(x) \, ds,
        \end{align*}
    where $L_{\widehat{w}^{(s)}}$ is the Laplacian with edge weights
        \begin{align*}
            \widehat{w}^{(s)}_{xy} = w_{xy}\, \theta\left(\frac{\rho_s(x)}{\pi(x)}, \,\frac{\rho_s(y)}{\pi(y)}\right).
        \end{align*}
    By the Cauchy--Schwarz inequality, with $\rho_0 = \mu$,
        \begin{align}\label{eq:cauchy-schwarz-1}
            |\rho_t(x) - \mu(x)|^2 &\leq \left|\int_0^t L_{\widehat{w}^{(s)}} \psi_s(x) \, ds \right|^2 \leq t \int_0^t |L_{\widehat{w}^{(s)}} \psi_s(x)|^2 \, ds \leq \int_0^1 |L_{\widehat{w}^{(s)}} \psi_s(x)|^2 \, ds.
        \end{align}
    Let $x\in V$ be fixed. Then, once again by Cauchy--Schwarz and the definition of $\widehat{w}^{(s)}_{xy}$, we have
        \begin{align*}
            |L_{\widehat{w}^{(s)}} \psi_s(x)|^2 &= \left|\sum_{y\sim x} \widehat{w}^{(s)}_{xy} (\psi_s(y) - \psi_s(x))\right|^2 \\
            &\leq  \left(\sum_{y\sim x} \widehat{w}^{(s)}_{xy} \right)\left(\sum_{y\sim x} \widehat{w}^{(s)}_{xy} (\psi_s(y) - \psi_s(x))^2 \right)\\
            &= \left(\sum_{y\sim x} w_{xy}\, \theta\left(\frac{\rho_s(x)}{\pi(x)}, \frac{\rho_s(y)}{\pi(y)}\right) \right)\left(\sum_{y\sim x} \widehat{w}^{(s)}_{xy} (\psi_s(y) - \psi_s(x))^2 \right).
        \end{align*}
    Define $r_z := \rho_s(z)/\pi(z)$ for $z\in V$. Using the bound $\theta(a,b)\le (a+b)/2$ for $a,b\ge 0$, we obtain
        \begin{align*}
            \sum_{y\sim x}\widehat{w}^{(s)}_{xy}
            &= \sum_{y\sim x} w_{xy}\,\theta(r_x,r_y)
            \le \frac{1}{2}r_x\sum_{y\sim x}w_{xy} + \frac{1}{2}\sum_{y\sim x}w_{xy}r_y.
        \end{align*}
    Since $\pi(x)=\Deg_x/\Vol_w(G)$, the first term satisfies
        \begin{align*}
            \frac{1}{2}r_x\sum_{y\sim x}w_{xy} = \frac{1}{2}r_x\Deg_x=\frac{1}{2}\Vol_w(G)\rho_s(x) \le \frac{1}{2}\Vol_w(G).
        \end{align*}
    For the second term, using $w_{xy}\le \Deg_y$, we have
        \begin{align*}
            \sum_{y\sim x}w_{xy}r_y
            &= \Vol_w(G)\sum_{y\sim x}\frac{w_{xy}}{\Deg_y}\rho_s(y)
            \le \Vol_w(G)\sum_{y\in V}\rho_s(y)=\Vol_w(G).
        \end{align*}
    Therefore, for each $x\in V$ and $0\le s\le 1$,
        \begin{align*}
            \sum_{y\sim x}\widehat{w}^{(s)}_{xy} \le \Vol_w(G).
        \end{align*}
    Plugging this into the bound for $|L_{\widehat{w}^{(s)}} \psi_s(x)|^2$, we have
        \begin{align}\label{eq:estimate-for-intL}
            \int_0^1 |L_{\widehat{w}^{(s)}} \psi_s(x)|^2\,ds &\leq \Vol_w(G)\int_0^1 \sum_{y\sim x} \widehat{w}^{(s)}_{xy} (\psi_s(y) - \psi_s(x))^2\,ds\notag\\
            &\leq \Vol_w(G)\int_0^1 \sum_{\{u,v\}\in E} \widehat{w}^{(s)}_{uv} (\psi_s(u) - \psi_s(v))^2\,ds\notag\\
            &= \Vol_w(G)\mathcal{A}((\psi_t), (\rho_t)).
        \end{align}
    Recalling~\cref{eq:cauchy-schwarz-1} and using~\cref{eq:estimate-for-intL}, we obtain
        \begin{align*}
            |\rho_t(x) - \mu(x)|^2 &\leq \int_0^1 |L_{\widehat{w}^{(s)}} \psi_s(x)|^2 \, ds \leq \Vol_w(G) \mathcal{A}((\psi_t), (\rho_t)) \leq \Vol_w(G)\left(\frac{2}{\lambda_2(L_{w^{(\mu)}})}\varepsilon^2 + \eta\right),
        \end{align*}
    and in particular it follows that
        \begin{align}\label{eq:bound-rho-mu}
            \frac{\rho_t(x)}{\mu(x)} &\leq 1 + \frac{1}{\mu(x)}\sqrt{\Vol_w(G)\left(\frac{2}{\lambda_2(L_{w^{(\mu)}})}\varepsilon^2 + \eta\right)}\notag\\
            &\leq 1 + \left(\min_{u\in V}\mu(u)\right)^{-1}\sqrt{\Vol_w(G)}\left(\sqrt{\frac{2}{\lambda_2(L_{w^{(\mu)}})}}\,|\varepsilon| + \sqrt{\eta}\right).
        \end{align}
    Equivalently, defining
        \begin{align*}
            C' &:= \left(\min_{u\in V}\mu(u)\right)^{-1}\sqrt{\frac{2\,\Vol_w(G)}{\lambda_2(L_{w^{(\mu)}})}},\\
            C''_{\eta} &:= \left(\min_{u\in V}\mu(u)\right)^{-1}\sqrt{\Vol_w(G)}\sqrt{\eta},
        \end{align*}
    we may write~\cref{eq:bound-rho-mu} in the form
        \begin{align*}
            \frac{\rho_t(x)}{\mu(x)} &\leq 1 + C' |\varepsilon| + C''_{\eta}.
        \end{align*}
    Thus, for each $\{x,y\}\in E$, we have
        \begin{align*}
            \theta\left(\frac{\rho_t(x)}{\pi(x)}, \frac{\rho_t(y)}{\pi(y)}\right)
            &= \theta\left(\frac{\mu(x)}{\pi(x)}\frac{\rho_t(x)}{\mu(x)},\,\frac{\mu(y)}{\pi(y)}\frac{\rho_t(y)}{\mu(y)}\right)\\
            &\leq \theta\left(\left(1+C'|\varepsilon|+C''_{\eta}\right)\frac{\mu(x)}{\pi(x)},\,\left(1+C'|\varepsilon|+C''_{\eta}\right)\frac{\mu(y)}{\pi(y)}\right)\\
            &= \left(1+C'|\varepsilon|+C''_{\eta}\right)\theta\left(\frac{\mu(x)}{\pi(x)},\,\frac{\mu(y)}{\pi(y)}\right),
        \end{align*}
    where the last step uses the fact that $\theta$ is 1-homogeneous. We thus have the bound, for each $f\in\mathbb{R}^{V}$ and $0\le t\le 1$,
        \begin{align*}
            f^\top L_{\widehat{w}^{(t)}} f &\leq \left(1 + C'|\varepsilon| + C''_{\eta}\right) f^\top L_{w^{(\mu)}} f,
        \end{align*} 
    and inversely so for $f\in\mathbb{R}^{V}_0$:
        \begin{align}\label{eq:bound-Z}
            f^\top L_{\widehat{w}^{(t)}}^{\dagger} f &\geq \frac{1}{1 + C'|\varepsilon| + C''_{\eta}} f^\top L_{w^{(\mu)}}^{\dagger} f.
        \end{align}
    Let $0\le t\le 1$ be fixed. The flux constraint requires that
        \begin{align*}
            \frac{d}{dt}\rho_t &= L_{\widehat{w}^{(t)}} \psi_t,
        \end{align*}
    and we may freely assume that $\psi_t\in\mathbb{R}^{V}_0$ since adding a constant to $\psi_t$ does not change the action or violate the flux constraint. Thus, since $L_{\widehat{w}^{(t)}}$ is positive definite on $\mathbb{R}^{V}_0$, $\psi_t$ is uniquely determined via
        \begin{align}\label{eq:psi-t-uniquely-determined}
            \psi_t &= L_{\widehat{w}^{(t)}}^{\dagger} \rho_{t}',
        \end{align}
    where $\rho_{t}' = \frac{d}{dt} \rho_t$ for brevity. We then have that, since $ L_{\widehat{w}^{(t)}}^{\dagger}  L_{\widehat{w}^{(t)}}  L_{\widehat{w}^{(t)}}^{\dagger}  =  L_{\widehat{w}^{(t)}}^{\dagger} $ and since $\rho_{t}'\in\mathbb{R}^{V}_0$ for all $t$, it holds
        \begin{align*}
            \mathcal{A}((\psi_t), (\rho_t)) &= \int_0^1 \psi_t^\top L_{\widehat{w}^{(t)}} \psi_t \, dt = \int_0^1 \left(\rho_{t}'\right)^\top L_{\widehat{w}^{(t)}}^{\dagger} \left(\rho_{t}'\right) dt \geq \frac{1}{1 + C'|\varepsilon| + C''_{\eta}} \int_0^1 (\rho_t')^\top L_{w^{(\mu)}}^\dagger (\rho_t') dt,
        \end{align*}
    where the last step follows from~\cref{eq:bound-Z}. Separately, by the Cauchy--Schwarz inequality,
        \begin{align*}
            \left(\int_0^1 \rho_t' \,dt\right)^\top L_{w^{(\mu)}}^{\dagger} \left(\int_0^1 \rho_t' \,dt\right) &= \left\|\int_0^1 (L_{w^{(\mu)}}^{\dagger})^{1/2} \rho_t'\, dt \right\|^2 \\
            &\leq \int_0^1 \|(L_{w^{(\mu)}}^{\dagger})^{1/2} \rho_t'\|^2 \,dt\\
            &= \int_0^1 (\rho_t')^\top L_{w^{(\mu)}}^\dagger (\rho_t') \,dt,
        \end{align*}
    so that
        \begin{align*}
            \frac{1}{1 + C'|\varepsilon| + C''_{\eta}} \int_0^1 (\rho_t')^\top L_{w^{(\mu)}}^\dagger (\rho_t') dt &\geq \frac{1}{1 + C'|\varepsilon| + C''_{\eta}} \left(\int_0^1 \rho_t' dt\right)^\top L_{w^{(\mu)}}^{\dagger} \left(\int_0^1 \rho_t' dt\right)\\
            &= \frac{1}{1 + C'|\varepsilon| + C''_{\eta}} (\rho_1 - \rho_0)^\top L_{w^{(\mu)}}^{\dagger} (\rho_1 - \rho_0) \\
            &= \frac{\varepsilon^2}{1 + C'|\varepsilon| + C''_{\eta}} h^\top L_{w^{(\mu)}}^{\dagger} h.
        \end{align*}
    Since $\Maas(\mu,\mu+\varepsilon h)^2 \leq \mathcal{A}((\psi_t),(\rho_t)) \leq \Maas(\mu,\mu+\varepsilon h)^2 + \eta$, the previous inequality yields
        \begin{align*}
            \Maas(\mu, \mu+\varepsilon \, h)^2 + \eta &\geq \frac{\varepsilon^2}{1 + C'|\varepsilon| + C''_{\eta}} h^\top L_{w^{(\mu)}}^{\dagger} h.
        \end{align*}
    Letting $\eta\to 0$ gives the desired lower bound
        \begin{align*}
            \Maas(\mu, \mu+\varepsilon \, h)^2 &\geq \frac{\varepsilon^2}{1 + C'|\varepsilon|} h^\top L_{w^{(\mu)}}^{\dagger} h,
        \end{align*}
    and the theorem is proved.
\end{proof}

\end{document}
